\subsection{Sifat-Sifat Perkalian Matriks}

Seperti telah dinyatakan di atas, kadang-kadang matriks dapat
dikalikan dalam satu urutan,
tetapi tidak dalam urutan lainnya. Namun, bahkan jika $AB$ dan $BA$ keduanya terdefinisi,
keduanya belum tentu sama. 

\begin{example}{Perkalian Matriks Tidak Komutatif}\label{exa:noncommutativemultiplication}
Bandingkan hasil kali $AB$ dan $BA$ untuk matriks $ A = \leftB
\begin{array}{rr}
1 & 2 \\
3 & 4
\end{array}
\rightB, B= \leftB
\begin{array}{rr}
0 & 1 \\
1 & 0
\end{array}
\rightB $ 
\end{example}

\begin{solution} 
Pertama, perhatikan bahwa $A$ dan $B$ sama-sama berukuran $2 \times 2$. Oleh karena itu, kedua
hasil kali $AB$ dan $BA$ terdefinisi.
Hasil kali pertama, yaitu $AB$, adalah
\begin{equation*}
AB = \leftB
\begin{array}{rr}
1 & 2 \\
3 & 4
\end{array}
\rightB \leftB
\begin{array}{rr}
0 & 1 \\
1 & 0
\end{array}
\rightB = \leftB
\begin{array}{rr}
2 & 1 \\
4 & 3
\end{array}
\rightB 
\end{equation*}
Hasil kali kedua, yaitu $BA$, adalah
\begin{equation*}
\leftB
\begin{array}{rr}
0 & 1 \\
1 & 0
\end{array}
\rightB \leftB
\begin{array}{rr}
1 & 2 \\
3 & 4
\end{array}
\rightB = \leftB
\begin{array}{rr}
3 & 4 \\
1 & 2
\end{array}
\rightB 
\end{equation*}
Oleh karena itu, $AB \neq BA$. 
\end{solution}

Contoh ini menunjukkan bahwa Anda tidak dapat mengasumsikan $AB=BA$ meskipun
perkalian terdefinisi dalam kedua urutan. Jika untuk suatu matriks $A$ dan
$B$ berlaku $AB=BA$, kita mengatakan bahwa $A$ dan $B$ \textbf{berkomutasi}\index{matriks!komutatif}. Ini merupakan salah satu
sifat penting perkalian matriks.

Berikut adalah sifat-sifat penting lainnya dari perkalian matriks.
Perhatikan bahwa sifat-sifat ini hanya berlaku ketika ukuran matriks-matriks sedemikian rupa sehingga hasil kalinya terdefinisi. 

\begin{proposition}{Sifat-Sifat Perkalian Matriks}\label{prop:propertiesofmatrixmultiplication}
Hal-hal berikut berlaku \index{perkalian matriks!sifat} untuk matriks $A,B,$ dan $C$, serta untuk skalar $r$ dan $s$,

\begin{equation}
A\left( rB+sC\right) =r\left( AB\right) +s\left( AC\right)  \label{matrixproperties1}
\end{equation}

\begin{equation}
\left( B+C\right) A=BA+CA  \label{matrixproperties2}
\end{equation}

\begin{equation}
A\left( BC\right) =\left( AB\right) C  \label{matrixproperties3}
\end{equation}
\end{proposition}

\ifdefined\showproofs
\begin{proof}
 Pertama, kita akan membuktikan \ref{matrixproperties1}. Kita akan menggunakan Definisi \ref{def:ijentryofproduct}
dan membuktikan pernyataan ini menggunakan entri pada posisi $(i,j)$ dari suatu matriks.
Oleh karena itu,
\begin{equation*}
\left( A\left( rB+sC\right) \right) _{ij}=\sum_{k}a_{ik}\left( rB+sC\right)
_{kj}=\sum_{k}a_{ik}\left( rb_{kj}+sc_{kj}\right)
\end{equation*}
\begin{equation*}
=r\sum_{k}a_{ik}b_{kj}+s\sum_{k}a_{ik}c_{kj}=r\left( AB\right) _{ij}+s\left(
AC\right) _{ij}
\end{equation*}
\begin{equation*}
=\left( r\left( AB\right) +s\left( AC\right) \right) _{ij}
\end{equation*}
Jadi, $A\left( rB+sC\right) =r(AB)+s(AC)$ seperti yang dinyatakan. 

Bukti \ref{matrixproperties2} mengikuti pola yang sama dan diserahkan sebagai latihan. 

Pernyataan \ref{matrixproperties3} adalah hukum asosiatif perkalian. Dengan menggunakan
Definisi \ref{def:ijentryofproduct},
\begin{equation*}
\left( A\left( BC\right) \right) _{ij}=\sum_{k}a_{ik}\left( BC\right)
_{kj}=\sum_{k}a_{ik}\sum_{l}b_{kl}c_{lj}
\end{equation*}
\begin{equation*}
=\sum_{l}\left( AB\right) _{il}c_{lj}=\left( \left( AB\right) C\right) _{ij}.
\end{equation*}
Ini membuktikan \ref{matrixproperties3}.
\end{proof}
\fi
